%% kapa: exam=1 round=21 year=2010 period=1 form=A subject=civil_law subject_ko=민법 seq=39 number=39
%% source: Q-Net 감정평가사 기출문제 (cst00309) artlSeq=5024181 file=extracted/1008/5024181/2010년 제21회 감정평가사 1차시험 기출문제/1교시.hwp
%% conversion: kordoc 4.24.51 -> tools/gen1.py (auto); review=proofread-sonnet; flags=none
\begin{question}{39}
소멸시효의 중단에 관한 설명으로 옳지 않은 것은? (다툼이 있으면 판례에 의함)
\begin{choices}
  \item 채무자가 소멸시효 완성 후에 채권자에 대하여 채무를 승인함으로써 그 시효의 이익을 포기한 경우에는 그때부터 새로이 소멸시효가 진행한다.
  \item 가압류에 의한 시효중단의 효력은 가압류의 집행보전의 효력이 존속하는 동안에는 계속된다.
  \item 시효의 중단은 당사자 및 그 승계인 간에만 효력이 있다.
  \item 물상보증인이 제기한 저당권설정등기 말소청구소송에서 저당권자가 응소하여 피담보채권의 존재를 적극 주장하였다면 피담보채권에 관한 소멸시효는 중단된다.
  \item 소유권이전등기를 명한 확정판결의 피고가 여전히 자신에게 소유권이 있다고 주장하면서 제기한 재심의 소도 시효중단사유가 된다.
\end{choices}
\end{question}
\answer{4}
